\scanpage{12}
\mspage{- 2 -}
n\textsuperscript{o} 398

PROPOSITION 2.- \uline{Si} $E'\longrightarrow E\longrightarrow E''$ \uline{est une suite exacte de} $A$-\uline{modules de longueur finie}, \uline{la suite} $E^{\prime\prime\ast}\longrightarrow E^\ast\longrightarrow E^{\prime\ast}$ \uline{est exacte}.

Soit d'abord $0\longrightarrow L\longrightarrow M\longrightarrow N\longrightarrow0$ une suite exacte de modules de longueur finie et montrons que $0\longrightarrow N^\ast\longrightarrow M^\ast\longrightarrow L^\ast$ est exacte. C'est en tout cas vrai pour $0\longrightarrow N^\ast\longrightarrow M^\ast\longrightarrow L^\ast$, d'après les propriétés générales de Hom. Comptant les longueurs, on en déduit que $M^\ast\longrightarrow L^\ast$ est surjectif, d'où le résultat.

Le cas d'une suite exacte qqe $E'\longrightarrow E\longrightarrow E''$ se ramène au précédent en dévissant à la façon habituelle, (Chap. II, 3e éd., p. 21). Soit $E$ un $A$-module. Si $F\subset E$, on note $F^{\unsure{o|\circ|0}}$ son orthogonal dans $E^\ast$, i.e. etc \ldots

PROPOSITION 3.- \uline{Soit} $E$ \uline{un} $A$-\uline{module de longueur finie}. \uline{La correspondance qui à} $F\subset E$ \uline{fait correspondre} $F^o$ \uline{est une bijection de l'ensemble des sous-modules de} $E$ \uline{sur celui des sous-modules de} $E^\ast$ ; \uline{cette bijection transforme somme en intersection et intersection en somme}.

On utilise la suite exacte $0\longrightarrow(E/F)^\ast\longrightarrow E^\ast\longrightarrow F^\ast\longrightarrow0$, qui montre que $F^\ast$ s'identifie à $E^\ast/F^o$ et $(E/F)^\ast$ à $F^o$. On en déduit que l'injection $F\longrightarrow F^{oo}$ s'identifie à l'injection $F\longrightarrow F^{\ast\ast}$, qui est un isomorphisme. D'où, etc \ldots

\uline{Exemple}.- $A={}$\uuline{$Z$}, $K/A={}$\uuline{$Q$}$/$\uuline{$Z$} ; on a une dualité pour les groupes abéliens finis.

\centerline{\texttt{-:-:-:-}}
% archive sheet number: 11

